\chapter{Среда исполнения}

Среда исполнения это плоскость, где живёт код и где происходят все эффекты. Она единственная в логическом смысле и пассивная в поведенческом: не запускает агентских циклов, не принимает решений, исполняет только то, что подтверждено журналом. Эта глава описывает, из чего она состоит и почему каждая часть устроена именно так.

\section{Авторитетная среда и реплики}

Логически состояние проекта одно: основная ветка и журнал определяют его полностью. Физически среда может быть пулом машин. Одна из них авторитетна: только на ней продвигается голова основной ветки. Остальные, исполнительные реплики, разделяют рабочее дерево по снимкам и обслуживают санбоксы задач. Отдельная роль, верификатор, повторяет критические прогоны тестов независимо от той машины, где работал продьюсер.

\begin{whybox}
Единственная физическая машина была бы и узким местом, и единственной точкой отказа. Пул реплик снимает ограничение ёмкости, а верификатор даёт хотя бы частичную независимость тестов от машины, которую мы вынуждены считать доверенной. Тезис о дорогой среде при этом не ломается: дорого воспроизвести среду на каждого контрибьютора, а не на две-три машины кор-команды.
\end{whybox}

\section{Санбокс: копия рабочего дерева не является границей}

Санбокс задачи состоит из двух вещей. Первая: копия рабочего дерева на момент назначения, созданная копированием при записи. На macOS это клон файлов APFS, который создаётся за миллисекунды и не занимает места, пока файл не изменён; на Linux снимок btrfs или ZFS либо overlay. Вторая: политика ограничения на уровне операционной системы.

\begin{warnbox}[Worktree это не безопасность]
Первая редакция протокола описывала санбокс как отдельный каталог с ограничением путей. Этого недостаточно. Копия рабочего дерева изолирует версии файлов, но не изолирует процессы и их сведения, сокеты и локальные сервисы, домашний каталог, общие кеши сборки, сокет Docker, временные каталоги, открытые порты, ресурсы. И главное: скрипты сборки. Постинсталляционный хук в npm или build-скрипт на Rust исполняется с правами агента, и ограничение путей в инструменте shell ему не указ. Границей может быть только механизм ядра.
\end{warnbox}

Политика ограничения задаёт: запись только в каталог санбокса и его временный каталог; чтение системы и тулчейна; запрет на остальной домашний каталог и локальные сервисы; запрет сигналов и сведений о чужих процессах; сеть закрыта, кроме локального egress-прокси, в котором живёт список разрешённых адресов. Агент никогда не получает shell хоста: каждая команда стартует внутри политики, и выйти из неё нельзя, не эксплуатируя ядро.

\subsection{Механизмы по платформам}

\begin{center}
\small
\begin{tabular}{p{2.3cm} p{5.3cm} p{6.0cm}}
\toprule
Платформа & Механизм & Что обеспечивает и чего нет \\
\midrule
macOS & Seatbelt (sandbox-exec), профиль на задачу; клон APFS; группа процессов со сторожем & файловая система, сеть, сигналы, сервисы: да; агрегатные лимиты памяти: нет, заменяет сторож; ядро общее \\
Linux, тысячи санбоксов & bubblewrap с user namespaces, seccomp, cgroups; overlay & всё, включая ресурсы; ядро общее \\
Linux, чужие агенты & gVisor & перехват системных вызовов, меньшая поверхность ядра \\
Linux, прочность VM & Firecracker, microVM на задачу & отдельное ядро; тысячи на хост штатно \\
macOS, прочность VM & macOS-гость в Virtualization & не больше двух гостей на хост; только для критических операций \\
\bottomrule
\end{tabular}
\end{center}

Seatbelt на macOS это тот же механизм, на котором стоит App Sandbox, и тот же, который используют собственные санбоксы вендоров агентских инструментов. Он даёт нулевую стоимость на процесс и тысячи одновременных профилей. Его ограничения честно перечислены в инвентаре (раздел 4.5): общее ядро и отсутствие агрегатных лимитов памяти.

\section{Слоты и очередь}

Санбокс дёшев: это каталог и политика, и в покое он не стоит ничего. Исполнение дорого: одновременно живых сборок на машине столько, сколько тянут процессор и память, для Xcode это десятки. Поэтому протокол различает санбоксы, которых могут быть тысячи, и слоты, число которых задаёт ёмкость среды. Назначение без свободного слота ждёт в очереди.

\begin{whybox}
Так устроен любой serverless: миллионы функций, сотни одновременных запусков. Полезный размер роя ограничен ёмкостью среды, а не числом контрибьюторов, и честнее сделать очередь частью протокола, чем обнаружить её в проде.
\end{whybox}

Отдельно про Xcode. DerivedData и кеш пакетов живут внутри санбокса, с прогретой базой, которая клонируется вместе с деревом. Симуляторы это разделяемый сервис, его нельзя честно посадить в санбокс; прогон на симуляторах и устройствах выполняется операцией раннера (раздел 4.4), а не из shell агента.

\section{Защищённые пути и операции раннера}

Тесты являются валидатором результата, поэтому защищаются сильнее кода. Тесты, конфигурация сборки, правила валидаторов и манифест проекта смонтированы во все санбоксы только на чтение. Их изменение проходит через задачи привилегированного класса (глава 6), санбоксы которых открываются с другой политикой.

Всё, что требует секретов проекта или общих сервисов, недоступно из shell вовсе. Подпись артефактов, публикация, прогоны на симуляторах и устройствах, доступ к данным предоставляются как узкие именованные операции доверенного раннера с очередью. Каждая операция авторизуется собственной записью журнала, имеет идентификатор и ключ идемпотентности и оставляет запись исхода: выполнена, отклонена или не установлен. Семантика повторов и поведение при смене аренды и остановке описаны в главе 7. Секреты проекта существуют только в раннере, секреты контрибьюторов не существуют на среде вовсе: это две разные половины инварианта I6. Агент вызывает операцию так же, как команду, но исполняет её не его санбокс, а раннер, у которого есть ключи и которому нельзя передать произвольную строку.

\begin{exbox}[Операции раннера в сценарии]
\texttt{sign\_artifact(hash)} подписывает сборку ключом проекта и возвращает хеш подписанного артефакта. \texttt{run\_device\_tests(candidate, suite)} ставит прогон в очередь на устройства и возвращает аттестацию. \texttt{publish(artifact, channel)} публикует. Ни в одной из них нет shell, и ни одна не принимает путь вне санбокса задачи.
\end{exbox}

\section{Receipts и инвентарь ограничений}

Каждая команда, которую агент отправляет в среду, получает идентификатор, digest самой команды и digest входного состояния санбокса, а также лимиты по времени, памяти, числу процессов и объёму вывода. По завершении среда выдаёт receipt: код выхода, дайджесты потоков вывода, длительность, потреблённые ресурсы, digest результирующего состояния. Receipts одной задачи связаны в хеш-цепочку, и голова цепочки входит в запись о ревизии (инвариант I10). Receipt удостоверяет происхождение, а не воспроизводимость и не полноту: он говорит, что среда получила такую-то команду и выдала такой-то результат, но ничего не говорит о том, что агент делал на своей машине вне протокола. Воспроизводимость требует отдельно сохранённых входов, digest окружения и тулчейна, зафиксированных зависимостей и, там, где они влияют на результат, внешних ответов; egress-прокси есть естественная точка, где внешние ответы фиксируются для повтора.

Модель команды и результата взята из API удалённого исполнения сборок: действие адресуется по содержимому, результат аттестован, повтор того же действия над тем же входом идемпотентен. Это позволяет верификатору повторить критическую команду и сравнить результат.

Инвентарь ограничений есть часть манифеста проекта: для каждого класса ограничения среда объявляет, обеспечивает она его полностью, частично или не обеспечивает. Задача объявляет требуемый класс. Несовпадение есть отказ, а не предупреждение (инвариант I9). На macOS инвентарь честно говорит: файловая система и сеть полностью, память частично, через сторож процесса.

\section{Светлый клиент журнала}

Среда исполняет только то, что подтверждено журналом (инвариант I1). Три точки входа:

\begin{itemize}
\item санбокс открывается по записи назначения с сертификатом кворума; каждая последующая команда несёт идентификатор задачи, подпись ноды и digest плана, и среда сверяет digest с записью;
\item операция раннера проводится по записи авторизации операции;
\item голова основной ветки продвигается по записи об интеграции с сертификатом кворума, атомарно, с проверкой, что текущая голова равна ожидаемой.
\end{itemize}

Среде не нужен весь журнал. Она хранит последний подписанный снапшот и проверяет включение конкретной записи по доказательству Меркла. Она хранит наивысшую принятую версию журнала как долговременное состояние безопасности и не принимает более раннюю (инвариант I8). И она проверяет актуальность разрешения на границе исполнения: запись, совершённая в прошлой эпохе или до остановки, не даёт права на исполнение сейчас, каким бы валидным ни был её сертификат. Компрометацию самой среды протокол не предотвращает; глава 9 разбирает, что именно это значит.

\section{Вход в санбокс}

Как человек с улицы попадает в санбокс на машине проекта и почему это не страшно. По шагам.

\begin{enumerate}
\item Идентичность без аккаунтов. У контрибьютора есть ключ ноды. Он делает join по адресу манифеста проекта, отправляет публичный ключ и объявленные оси, получает план и разницу полномочий. Кор-нода записывает ноду в журнал. Ни паролей, ни регистрации: запись о ноде есть, прав нет.
\item Права появляются только с назначением. Оркестратор создаёт запись назначения: ключ ноды, задача, digest плана, класс санбокса из контракта (платформа, процессор, память, диск, сеть, время) и слот. До совершения этой записи среда о контрибьюторе ничего не знает.
\item Санбокс заводится по записи, а не по запросу. Демон видит совершённое назначение через доказательство включения и провизионирует санбокс нужного класса с лимитами из контракта и egress-прокси. Санбокс живёт от назначения до закрытия задачи.
\item Аутентификация это сертификат на задачу. Среда держит свой CA. При совершении назначения демон выпускает короткоживущий SSH-сертификат для транспортного ключа ноды: principal равен идентификатору задачи, срок равен окну назначения, принудительная команда запускает демон по стандартным потокам ровно для этого санбокса. Сертификат приходит в ответ на назначение. Контрибьютор подключается своим ключом с этим сертификатом и оказывается сразу внутри санбокса. Отзыв не нужен: сертификат истекает сам, а закрытие задачи или остановка рвут сессии.
\item Каждая команда несёт digest плана, среда сверяет его с записью и выдаёт receipt.
\item Сетевой путь. Среда за NAT в офисе публикуется через маленький relay на VPS, который пробрасывает SSH в tailnet среды; среда на публичном адресе принимает только сертификаты, без паролей и без ключей в authorized\_keys.
\end{enumerate}

\begin{warnbox}[{SSH сразу в санбокс: безопасно при пяти условиях}]
Паттерн стандартный: так работают git-хостинги с принудительной командой, облачные среды разработки с SSH в контейнер, отладочные сессии CI. Небезопасным он становится, когда через SSH можно получить что-то, кроме санбокса. Условия. Первое: в сертификате опция restrict, никакого перенаправления портов, agent forwarding, X11, sftp, переменных окружения и pty; клиент не может запросить другую команду. Второе: команда исполняется от отдельного служебного пользователя без прав и без секретов; на Linux дальше user namespaces дают санбоксу свой uid, на macOS разделение процессов разных задач обеспечивает Seatbelt. Третье: только сертификаты от CA среды, без паролей и без ключей в файле; транспортный ключ сертифицирован ключом ноды и не совпадает с ключом подписи ревизий, один ключ на два протокола плохая практика. Четвёртое: поверхность sshd. OpenSSH работает от root, и уязвимость до аутентификации означает компрометацию машины; для публичного роя SSH-сервер встраивается в сам непривилегированный демон на библиотеке и принимает только exec протокола по stdio, root в приёме соединений не участвует; для L0 и корпоративного варианта достаточно OpenSSH с первым и третьим условиями. Пятое: границей всё равно остаётся санбокс. SSH это аутентификация, связывающая сессию с назначением, а не изоляция; если санбокс слаб, безопасный SSH не спасёт.
\end{warnbox}

Что защищает от человека с улицы, не аутентификация. Защищают граница ядра в санбоксе, отсутствие секретов на среде, allowlist сети, лимиты и слоты, и то, что в проект ничего не попадает без ревью другой ноды. Доказательства идентичности (глава 8) здесь необязательны: они меняют вес ноды и глубину ревью, а не право войти в свой санбокс.

\section{Shell, неотличимый от локального}

Удалённый shell есть единственный инструмент рабочей ноды, поэтому его качество определяет качество работы всего роя. Требование сформулировано так: для модели shell среды неотличим от локального, не в смысле обмана, а в смысле того, что замечать нечего.

Механика, которая это даёт, называется восстановленной сессией. Каждый вызов есть отдельный процесс: точный код выхода, раздельные потоки вывода, отмена по группе процессов без последствий для следующего вызова. При этом состояние сессии, каталог, переменные, функции, снимается после команды по отдельному дескриптору, мимо потоков вывода, и восстанавливается перед следующей командой. Модель делает \texttt{cd src} в одном вызове и оказывается там в следующем, как в терминале.

Мутирующие вызовы одной сессии сериализованы: две команды, одновременно меняющие каталог, не могут потерять изменение друг друга, они просто выполняются по очереди. Параллельность даётся именованными сессиями, как вкладками терминала. Длительные процессы не переживают вызов, поэтому они запускаются явно в фоне и живут по дескриптору. Вывод сверх лимита остаётся в санбоксе, модель получает начало и конец и дальше читает его тем же shell.

\begin{warnbox}
Восстановленная сессия есть механизм качества, а не безопасности. Команда модели может испортить снимок своей сессии; она не получает от этого ничего, чего у неё не было. Границей остаётся санбокс.
\end{warnbox}

Транспорт: SSH только для аутентификации по сертификату задачи (раздел 4.7), а поверх него собственный протокол по стандартным потокам ввода-вывода демона среды. Ни отдельной аутентификации, ни shell хоста, один тёплый канал с мультиплексированием вызовов. Демон среды принадлежит проекту, его хеш опубликован в манифесте, и клиент контрибьютора пинит ожидаемую идентичность. Тот же протокол поверх другого транспорта превращает среду в облаке или в контейнере в ту же самую среду с точки зрения агента.
